\documentclass[12pt, a4paper]{article}
\usepackage[top=1cm, bottom=1.5cm, left=1cm, right=1cm]{geometry}
\usepackage{amsmath,amssymb,amsthm,mathtools}
\usepackage{graphicx}
\usepackage{braket}
\title{Devoir de BSM du 19 Mai 2026} 
\author{
	RAZAFIMAHERY Faneva Liantsoa\\
	\medskip % Un petit espace
	\normalsize Master 2 -Physique des Hautes Energies \\
	\medskip
	\normalsize Université d'Antananarivo  \\
	%\vspace{1cm} %  pour séparer verticalement bro
	% Ensuite, on insère le titre manuellement
	%\Large \textbf{DEVOIR QFT: Théorème de Wick} \\ si on veut mettre le titre ici et en gras
	% \Huge C'est grand et \LARGE ou \Large c plus petit
}
\date{30 Mai 2026} 
\begin{document}
	
	\maketitle
	\begin{abstract}
		Dans ce travail, nous allons montrer quelques relations dans le cadre du formalisme d'hélicité spinorielle en partant des définitions déja connues suivantes:
		
		\begin{equation}
			\begin{split}
				 (i)\quad &p_{a \dot{a}}=\begin{pmatrix}
					p_0 + p_3 & p_1-ip_2\\
					p_1 + i_2 & p_0 - p_3
				\end{pmatrix}\\
				 (ii)\quad&\mathbf{det} (p_{a \dot{a}})=m^2\\
				 (iii)\quad&\bar{p}^{a \dot{a}}=(\varepsilon^{ab} p_{b \dot{b}} \varepsilon^{\dot{b}\dot{a}})^{\mathbf{T}}\\
				 (iv)\quad& \varepsilon^{a b}=\begin{pmatrix}
					0 & 1\\
					-1 & 0
				\end{pmatrix}\\
				 (v)\quad& \varepsilon^{\dot{a}\dot{b}}=\begin{pmatrix}
					0 & -1\\
					1 & 0
				\end{pmatrix}\\
				(vi)\quad& \mathbf{tr}(p_{a \dot{a}}\bar{q}^{\dot{a} b})= 2p_{\mu}q^{\mu}
			\end{split}
		\end{equation}
		
	\end{abstract} 
	\section{}
		\subsection{Soit à montrer que si \(m=0\),\quad on a \(p_{a \dot{a}}=\lambda_a \bar{\lambda}_{\dot{a}}\)}
		
		En utilisant la relation \((ii)\), on trouve que:
		\begin{equation}
			m=0 \Leftrightarrow \mathbf{det}(p_{a \dot{a}})=0
		\end{equation}
		Grâce à l'algèbre linéaire, nous savons que pour une matrice \(2\times2\), on a:
		\begin{equation}
		\mathbf{det}(p_{a \dot{a}})=0 \Leftrightarrow \mathbf{Rank}(p_{a \dot{a}}) < 2
		\end{equation}
		Or une matrice de rang \(0\) aurait tous ses éléments égaux à \(0\).
		Ce n'est pas le cas de la matrice \(p_{a \dot{a}}\).\\
		On en déduit alors que \(\mathbf{Rank}(p_{a \dot{a}})=1\).\\
		D'après l'algèbre linéaire, on peut alors écrire:
		\begin{equation}
			p_{a \dot{a}}=\lambda_a \bar{\lambda}_{\dot{a}} \quad\text{avec}\begin{cases}
				&\lambda_a  \quad \text{un vecteur colonne}\\
				&\lambda_{\dot{a}} \quad \text{un vecteur ligne}\\
			\end{cases} 
		\end{equation}
		Nous avons donc montré la relation demandée.
		
		\subsection{Soit à en déduire \(\lambda_1\), \(\lambda_2\), \(\lambda_{\dot{1}}\), \(\lambda_{\dot{2}}\)}
		
		Procédons directement en écrivant la décomposition suivante
	\begin{equation}
		p_{a \dot{a}}=
		\begin{pmatrix}
			k\\
			l
		\end{pmatrix}  \times 
		\begin{pmatrix}
			m & n
		\end{pmatrix}	
	\end{equation}
	En utilisant la relation \((i)\), on peut écrire que:
	\begin{equation}
			p_{a \dot{a}}=
			\begin{pmatrix}
				km & kn\\
				lm & ln
			\end{pmatrix} = 
			\begin{pmatrix}
			p_0 + p_3 & p_1-ip_2\\
			p_1 + i_2 & p_0 - p_3
			\end{pmatrix} 
	\end{equation}
	De cette égalité, nous tirons les égalités suivantes:
	\begin{equation}
		\begin{split}
			&km= p_0 + p_3\\
			&k=\alpha \sqrt{p_0 + p_3}\\
			&m=\frac{1}{\alpha} \sqrt{p_0 + p_3}\\
			&kn=p_1-ip_2 \Rightarrow n= \frac{p_1-ip_2}{k}=\frac{1}{\alpha} \frac{p_1-ip_2}{ \sqrt{p_0 + p_3}}\\
			&lm=p_1 + i_2 \Rightarrow l= \frac{p_1+ip_2}{m}= 
			\alpha \frac{p_1+ip_2}{ \sqrt{p_0 + p_3}}\\
		\end{split}
	\end{equation}
	En posant \(\alpha=1\), on trouve que:
	\begin{equation}
	p_{a \dot{a}}=
	\begin{pmatrix}
		k\\
		l
	\end{pmatrix}  \times 
	\begin{pmatrix}
		m & n
	\end{pmatrix}	=
	\begin{pmatrix}
		\sqrt{p_0 + p_3}\\
		\frac{p_1+ip_2}{ \sqrt{p_0 + p_3}}
	\end{pmatrix}  \times 
	\begin{pmatrix}
		\sqrt{p_0 + p_3} & \frac{p_1-ip_2}{ \sqrt{p_0 + p_3}}
	\end{pmatrix}
	\end{equation}
	Or, on avait \(p_{a \dot{a}}=\lambda_a \lambda_{\dot{a}}\), ce qui nous donne donc que :
	\begin{equation}
		\begin{split}
			\lambda_a &= \begin{pmatrix}
				\sqrt{p_0 + p_3}\\
				\frac{p_1+ip_2}{ \sqrt{p_0 + p_3}}
			\end{pmatrix}\\
			\lambda_{\dot{a}} &= \begin{pmatrix}
				\sqrt{p_0 + p_3} & \frac{p_1-ip_2}{ \sqrt{p_0 + p_3}}
			\end{pmatrix}\\
		\end{split}
	\end{equation}
	Par identification des indices des spineurs, on trouve que:
	
	\begin{center}\fbox{%
		\begin{minipage}{0.4\textwidth}
			\begin{equation}
				\begin{split}
				\lambda_1&=\sqrt{p_0 + p_3}\\
				\lambda_2&= \frac{p_1+ip_2}{ \sqrt{p_0 + p_3}}\\
				\lambda_{\dot{1}}&=\sqrt{p_0 + p_3}\\
				\lambda_{\dot{2}}&=\frac{p_1-ip_2}{ \sqrt{p_0 + p_3}}
				\end{split}
			\end{equation}%
		\end{minipage}
		}
	\end{center}
	\section{Soit à montrer que \(p_{\mu}q^{\mu}=\frac{1}{2}(\varepsilon^{ab}\lambda_a \mu_b)(\varepsilon^{\dot{a}\dot{b}}\bar{\lambda}_{\dot{a}} \bar{\mu}_{\dot{b}})\)}
	
	On a: \(p.q= g_{\mu\nu}p^{\mu}q^{\nu}\)\\
	Et de la définition \(p_{a \dot{a}}=\sigma_{a \dot{a}}^{\mu}p_{\mu}\), on peut tirer:
	\begin{equation}
		p^{\mu}=\frac{1}{2}\sigma^{a \dot{a}\mu}p_{a \dot{a}}
	\end{equation}
	Ainsi, on peut écrire: 
	\begin{equation}
		\begin{split}
			p.q
			&= g_{\mu\nu}p^{\mu}q^{\nu}\\
			&= g_{\mu\nu} \frac{1}{4}\sigma^{a \dot{a}\mu}p_{a \dot{a}}\sigma^{b \dot{b}\nu}q_{b \dot{b}}
		\end{split}
	\end{equation}
	Or, on a: \(g_{\mu\nu}\sigma^{a \dot{a}\mu}\sigma^{b \dot{b}\nu}= 2\varepsilon^{a \dot{a}}\varepsilon^{b \dot{b}}\)\\
	Celà nous donne ainsi: 
	\begin{equation}
		\begin{split}
			p.q
			&= g_{\mu\nu} \frac{1}{4}\sigma^{a \dot{a}\mu}p_{a \dot{a}}\sigma^{b \dot{b}\nu}q_{b \dot{b}}\\
			&=\frac{1}{2}\varepsilon^{a \dot{a}}p_{a \dot{a}}\varepsilon^{b \dot{b}}q_{b \dot{b}}\\
		\end{split}
	\end{equation}
	Or, on pose \(m=0\), ce qui nous donne: \(p_{a \dot{a}}=\lambda_a \bar{\lambda}_{\dot{a}}\) et \(q_{b \dot{b}}=\mu_b \bar{\mu}_{\dot{b}}\).\\
	Finalement, on a bien la forme:
	
	\begin{center}
	\fbox{\(p_{\mu}q^{\mu}=\frac{1}{2}(\varepsilon^{ab}\lambda_a \mu_b)(\varepsilon^{\dot{a}\dot{b}}\bar{\lambda}_{\dot{a}} \bar{\mu}_{\dot{b}})\)}
	\end{center}
	
	
	
	
	\section{Soit à montrer que \(\lambda^1_a \braket{\lambda^2, \lambda^3} + \lambda^2_a \braket{\lambda^3, \lambda^1} + \lambda^3_a
	\braket{\lambda^1, \lambda^2}=0
	\)}
	On définit: 
	\begin{equation}
		\braket{\lambda^i, \lambda^j}= \varepsilon^{ab} \lambda^i_a \lambda^j_b
	\end{equation}
	Ainsi, on a:
	\begin{equation}
		\begin{split}
		\lambda^1_a \braket{\lambda^2, \lambda^3} + \lambda^2_a \braket{\lambda^3, \lambda^1} + \lambda^3_a
		\braket{\lambda^1, \lambda^2}
		=
		&\lambda^1_a \varepsilon^{bc} \lambda^2_b \lambda^3_c + 
		\lambda^2_a \varepsilon^{bc} \lambda^3_b \lambda^1_c +
		\lambda^3_a \varepsilon^{bc} \lambda^1_b \lambda^2_c \\
		=
		&\lambda^1_a \varepsilon^{12} \lambda^2_1 \lambda^3_2 +
		\lambda^1_a \varepsilon^{21} \lambda^2_2 \lambda^3_1 +
		\lambda^2_a \varepsilon^{12} \lambda^3_1 \lambda^1_2 +
		\lambda^2_a \varepsilon^{21} \lambda^3_2 \lambda^1_1 \\
		&+\lambda^3_a \varepsilon^{12} \lambda^1_1 \lambda^2_2 +
		\lambda^3_a \varepsilon^{21} \lambda^1_2 \lambda^2_1\\
		\end{split}
	\end{equation}
	Or, on sait que \(\varepsilon^{bc}=-\varepsilon^{cb}\).\\
	Ce qui nous donne 
	\begin{equation}
		\begin{split}
			\lambda^1_a \braket{\lambda^2, \lambda^3} + \lambda^2_a \braket{\lambda^3, \lambda^1} + \lambda^3_a
			\braket{\lambda^1, \lambda^2}
			=
			&\lambda^1_a \varepsilon^{12} \lambda^2_1 \lambda^3_2 -
			\lambda^1_a \varepsilon^{12} \lambda^2_2 \lambda^3_1 +
			\lambda^2_a \varepsilon^{12} \lambda^3_1 \lambda^1_2 -
			\lambda^2_a \varepsilon^{12} \lambda^3_2 \lambda^1_1 \\
			&+\lambda^3_a \varepsilon^{12} \lambda^1_1 \lambda^2_2 -
			\lambda^3_a \varepsilon^{12} \lambda^1_2 \lambda^2_1\\
		\end{split}
	\end{equation}
	\begin{itemize}
	\item Pour le cas \(a=1\), on a:
	\begin{equation}
		\begin{split}
			\lambda^1_1 \braket{\lambda^2, \lambda^3} + \lambda^2_1 \braket{\lambda^3, \lambda^1} + \lambda^3_1
			\braket{\lambda^1, \lambda^2}
			=
			&\lambda^1_1 \varepsilon^{12} \lambda^2_1 \lambda^3_2 -
			\lambda^1_1 \varepsilon^{12} \lambda^2_2 \lambda^3_1 +
			\lambda^2_1 \varepsilon^{12} \lambda^3_1 \lambda^1_2 -
			\lambda^2_1 \varepsilon^{12} \lambda^3_2 \lambda^1_1 \\
			&+\lambda^3_1 \varepsilon^{12} \lambda^1_1 \lambda^2_2 -
			\lambda^3_1 \varepsilon^{12} \lambda^1_2 \lambda^2_1\\
			=
			&\lambda^1_1 \varepsilon^{12} \lambda^2_1 \lambda^3_2 -
			\lambda^2_1 \varepsilon^{12} \lambda^3_2 \lambda^1_1 +
			\lambda^2_1 \varepsilon^{12} \lambda^3_1 \lambda^1_2 -
			\lambda^3_1 \varepsilon^{12} \lambda^1_2 \lambda^2_1 \\
			&+\lambda^3_1 \varepsilon^{12} \lambda^1_1 \lambda^2_2 -
			\lambda^1_1 \varepsilon^{12} \lambda^2_2 \lambda^3_1 \\
			=&0\\
		\end{split}
	\end{equation}
	\item Pour le cas \(a=2\), on a:
	\begin{equation}
		\begin{split}
			\lambda^1_2 \braket{\lambda^2, \lambda^3} + \lambda^2_2 \braket{\lambda^3, \lambda^1} + \lambda^3_2
			\braket{\lambda^1, \lambda^2}
			=
			&\lambda^1_2 \varepsilon^{12} \lambda^2_1 \lambda^3_2 -
			\lambda^1_2 \varepsilon^{12} \lambda^2_2 \lambda^3_1 +
			\lambda^2_2 \varepsilon^{12} \lambda^3_1 \lambda^1_2 -
			\lambda^2_2 \varepsilon^{12} \lambda^3_2 \lambda^1_1 \\
			&+\lambda^3_2 \varepsilon^{12} \lambda^1_1 \lambda^2_2 -
			\lambda^3_2 \varepsilon^{12} \lambda^1_2 \lambda^2_1\\
			=
			&\lambda^1_2 \varepsilon^{12} \lambda^2_1 \lambda^3_2 -
			\lambda^3_2 \varepsilon^{12} \lambda^1_2 \lambda^2_1 +
			\lambda^2_2 \varepsilon^{12} \lambda^3_1 \lambda^1_2 -
			\lambda^1_2 \varepsilon^{12} \lambda^2_2 \lambda^3_1  \\
			&+\lambda^3_2 \varepsilon^{12} \lambda^1_1 \lambda^2_2 -
			\lambda^2_2 \varepsilon^{12} \lambda^3_2 \lambda^1_1 \\
			=&0\\
		\end{split}
	\end{equation}
	\end{itemize}
	Ainsi on a:
	\fbox{\(\lambda^1_a \braket{\lambda^2, \lambda^3} + \lambda^2_a \braket{\lambda^3, \lambda^1} 		+ 	\lambda^3_a
		\braket{\lambda^1, \lambda^2}=0 \qquad \text{pour}\qquad a=1,2\)}
		
	
	
	
	
	
	
	

\end{document}